% !TEX program = xelatex
\documentclass[11pt,letterpaper]{article}

% ============================================================
% DATOS DEL CLIENTE  <-- cambiar aqui y recompilar
% ============================================================
\newcommand{\clienteCorto}{\PENDIENTE{cliente}{nombre corto}}
\newcommand{\clienteLargo}{\PENDIENTE{cliente}{razón social}}
\newcommand{\destinatario}{\PENDIENTE{cliente}{destinatario}}
\newcommand{\fechaProp}{\PENDIENTE{Fernando}{fecha}}
\newcommand{\demoURL}{pharmaweigh.appsoluciones.duckdns.org}
\newcommand{\tarifa}{\PENDIENTE{Fernando}{tarifa}}
\newcommand{\docTipo}{Resumen Ejecutivo}
\newcommand{\docPrefijo}{PROPUESTA}

% ============================================================
% maw_comun.tex --- preambulo compartido MAW Soluciones
% Lo usan PROPUESTA, CONTRATO y RESUMEN_EJECUTIVO.
% El documento debe definir ANTES del \input:
%   \docTipo, \docPrefijo, \fechaProp, \tarifa, \demoURL,
%   \clienteCorto, \clienteLargo
% Compilar con xelatex (Montserrat via fontspec), dos pasadas.
% ============================================================

% ============================================================
% PAQUETES
% ============================================================
\usepackage[spanish,es-tabla,es-noindentfirst]{babel}
\usepackage{fontspec}
\usepackage{montserrat}
\renewcommand*\familydefault{\sfdefault}
\usepackage[T1]{fontenc}
\usepackage{geometry}
\usepackage{graphicx}
\usepackage[table]{xcolor}
\usepackage{tabularx}
\usepackage{array}
\usepackage{booktabs}
\usepackage{ragged2e}
\usepackage{enumitem}
\usepackage{tcolorbox}
\usepackage{needspace}
\usepackage{fancyhdr}
\usepackage{siunitx}
\usepackage{hyperref}
\usepackage{tikz}
\usetikzlibrary{positioning, arrows.meta, shapes.geometric, calc}

\geometry{top=2cm, bottom=2cm, left=1.8cm, right=1.8cm}
\setlength{\headheight}{26pt}
\addtolength{\topmargin}{-14pt}

\sisetup{group-separator={,}, group-minimum-digits=4, output-decimal-marker={.}}

% ============================================================
% COLORES --- MAW Soluciones / PharmaWeigh
% (los mismos tokens que usa la app: src/app/globals.css)
% ============================================================
\definecolor{azulcom}{HTML}{1E40AF}    % acento principal (primary)
\definecolor{doradocom}{HTML}{2563EB}   % filete y acento secundario (blue-600)
\definecolor{grisbar}{HTML}{1E293B}     % slate-800: barra lateral de la app
\definecolor{grisfondo}{HTML}{F1F5F9}   % slate-100: fondo de tablas
\definecolor{gristexto}{HTML}{64748B}   % slate-500: texto secundario
\definecolor{verdeok}{HTML}{16A34A}     % éxito / dentro de tolerancia
\definecolor{ambarpend}{HTML}{D97706}   % advertencia / pendiente
\definecolor{rojonota}{HTML}{DC2626}    % error / fuera de tolerancia
\definecolor{tinta}{HTML}{0F172A}       % texto principal

\hypersetup{colorlinks=true, linkcolor=azulcom, urlcolor=doradocom}

% ============================================================
% ENCABEZADO / PIE
% ============================================================
\pagestyle{fancy}
\fancyhf{}
\fancyhead[L]{\footnotesize\color{gristexto}\textbf{MAW Soluciones} \textbar{} \docTipo}
\fancyhead[R]{\footnotesize\color{gristexto}\fechaProp}
\fancyfoot[C]{\footnotesize\color{gristexto}\thepage}
\renewcommand{\headrulewidth}{0.6pt}
\renewcommand{\headrule}{\hbox to\headwidth{\color{azulcom}\leaders\hrule height \headrulewidth\hfill}}

% ============================================================
% TITULOS Y CAJAS
% ============================================================
\newcommand{\tituloCot}[1]{%
    \noindent{\LARGE\bfseries\color{azulcom}\docPrefijo{} \textbar{} #1}\par
    \vspace{2pt}
    {\color{doradocom}\rule{\linewidth}{1.8pt}}\par
    \vspace{14pt}
}

\newcommand{\seccion}[1]{%
    {\large\bfseries\color{azulcom}#1}\\[4pt]
    {\color{doradocom}\rule{\linewidth}{1pt}}\par
    \vspace{8pt}
}

\newcommand{\subseccion}[1]{%
    \vspace{4pt}
    \noindent{\bfseries\color{azulcom}#1}\par
    \vspace{4pt}
}

\newtcolorbox{cajaNota}{
    colback=grisfondo, colframe=gristexto!60, rounded corners,
    boxrule=1pt, left=10pt, right=10pt, top=6pt, bottom=6pt
}

\newtcolorbox{cajaInfo}[1][]{
    colback=azulcom!5, colframe=azulcom,
    fonttitle=\bfseries\color{white}, colbacktitle=azulcom, title=#1,
    rounded corners, boxrule=1pt, left=8pt, right=8pt, top=6pt, bottom=6pt
}

\newtcolorbox{cajaTotal}{
    colback=verdeok!8, colframe=verdeok, rounded corners,
    boxrule=1.5pt, left=10pt, right=10pt, top=10pt, bottom=10pt
}

\newtcolorbox{cajaAviso}[1][]{
    colback=rojonota!6, colframe=rojonota,
    fonttitle=\bfseries\color{white}, colbacktitle=rojonota, title=#1,
    rounded corners, boxrule=1.2pt, left=8pt, right=8pt, top=6pt, bottom=6pt
}

% ============================================================
% PENDIENTES --- nada que el cliente o Fernando no hayan
% confirmado se escribe como cifra. Se marca y se ve.
% ============================================================
\newtcolorbox{cajaPendiente}[1][Pendiente de autorización]{
    colback=ambarpend!8, colframe=ambarpend,
    fonttitle=\bfseries\color{white}, colbacktitle=ambarpend, title=#1,
    rounded corners, boxrule=1.2pt, left=8pt, right=8pt, top=6pt, bottom=6pt
}

% \PENDIENTE{quien}{que} --- marca visible en linea
\newcommand{\PENDIENTE}[2]{%
    {\setlength{\fboxsep}{2.5pt}%
     \colorbox{ambarpend!20}{\color{ambarpend!80!black}\footnotesize\bfseries PENDIENTE (#1): #2}}%
}

% Version compacta para celdas de tabla
\newcommand{\PENDmini}{%
    {\setlength{\fboxsep}{2pt}%
     \colorbox{ambarpend!20}{\color{ambarpend!80!black}\scriptsize\bfseries pend.}}%
}

% ============================================================
% MOTOR DE CIFRAS
% ------------------------------------------------------------
% Mientras \cifrasDefinidas sea falso, toda hora e importe se
% imprime como PENDIENTE. Para activar el documento:
%   1. Fernando autoriza tarifa y horas.
%   2. Se pone \cifrasDefinidastrue y se llenan \tarifaNum y
%      los \def de horas de cada renglon.
%   3. Los subtotales, el total, la mensualidad y el punto de
%      equilibrio se calculan solos con estas macros.
% ============================================================
\newif\ifcifrasDefinidas
\cifrasDefinidasfalse

\providecommand{\tarifaNum}{0}     % MXN por hora de ingenieria (sin separadores)

% \HR{expresion} --- horas (acepta sumas: \HR{\hAa+\hAb})
\newcommand{\HR}[1]{\ifcifrasDefinidas\fpeval{#1}\else\PENDmini\fi}

% \IM{expresion} --- importe = horas x tarifa
\newcommand{\IM}[1]{\ifcifrasDefinidas\$\num{\fpeval{round((#1)*\tarifaNum,0)}}\else\PENDmini\fi}

% \MXN{numero} --- importe literal ya autorizado (mensualidad, infra)
\newcommand{\MXN}[1]{\ifcifrasDefinidas\$\num{#1}\else\PENDmini\fi}

% \RAZON{a}{b} --- cociente redondeado hacia arriba (punto de equilibrio)
\newcommand{\RAZON}[2]{\ifcifrasDefinidas\num{\fpeval{ceil((#1)/(#2))}}\else\PENDmini\fi}

% ============================================================
% TABLAS
% ============================================================
\newcolumntype{R}[1]{>{\RaggedLeft\arraybackslash}p{#1}}
\newcolumntype{C}[1]{>{\centering\arraybackslash}p{#1}}
\newcolumntype{L}[1]{>{\RaggedRight\arraybackslash}p{#1}}

\renewcommand{\arraystretch}{1.3}


% ============================================================
% Mismas macros de cifras que la propuesta.
% ============================================================
% \cifrasDefinidastrue
\def\tarifaNum{0}

\def\hAa{0}\def\hAb{0}\def\hAc{0}\def\hAd{0}
\def\hBa{0}\def\hBb{0}\def\hBc{0}\def\hBd{0}
\def\hCa{0}\def\hCb{0}\def\hCc{0}
\def\hDa{0}\def\hDb{0}\def\hDc{0}
\def\hEa{0}\def\hEb{0}\def\hEc{0}\def\hEd{0}\def\hEe{0}
\def\hFa{0}\def\hFb{0}\def\hFc{0}
\def\hGa{0}\def\hGb{0}
\def\hHa{0}\def\hHb{0}
\def\hIa{0}\def\hIb{0}\def\hIc{0}\def\hId{0}
\def\hJa{0}\def\hJb{0}\def\hJc{0}
\def\hKa{0}\def\hKb{0}\def\hKc{0}\def\hKd{0}
\def\hMa{0}\def\hMb{0}\def\hMc{0}

\def\sA{\hAa+\hAb+\hAc+\hAd}
\def\sB{\hBa+\hBb+\hBc+\hBd}
\def\sC{\hCa+\hCb+\hCc}
\def\sD{\hDa+\hDb+\hDc}
\def\sE{\hEa+\hEb+\hEc+\hEd+\hEe}
\def\sF{\hFa+\hFb+\hFc}
\def\sG{\hGa+\hGb}
\def\sH{\hHa+\hHb}
\def\sI{\hIa+\hIb+\hIc+\hId}
\def\sJ{\hJa+\hJb+\hJc}
\def\sK{\hKa+\hKb+\hKc+\hKd}
\def\sTOTAL{\sA+\sB+\sC+\sD+\sE+\sF+\sG+\sH+\sI+\sJ+\sK}

\begin{document}

\tituloCot{PharmaWeigh}

\noindent\RaggedRight{\small\color{gristexto}Cliente: \textbf{\clienteLargo} \quad\textbar\quad Preparado por: \textbf{MAW Soluciones} \quad\textbar\quad Demo en vivo: \href{https://\demoURL}{\demoURL} \quad\textbar\quad Vigencia: 30 días naturales}

\vspace{4pt}

\noindent{\small\color{gristexto}Este documento resume la \textbf{Propuesta PharmaWeigh}. El desglose completo de horas, bloques y condiciones está en la propuesta; los términos legales, en el contrato.}

\vspace{6pt}

\noindent{\small\color{ambarpend!80!black}\textbf{Borrador interno:} las cifras y los datos del cliente van marcados en ámbar porque no están autorizados. El contenido funcional sí es firme.}

\vspace{6pt}

\noindent Estimado \destinatario:

\vspace{6pt}

% ------------------------------------------------------------
\seccion{El problema}

\noindent En el área de pesaje, el error caro no es el peso mal tomado: es el que nadie detiene a tiempo. Un lote equivocado, en cuarentena, caducado o un gramaje fuera de tolerancia se descubre cuando el producto ya está mezclado. El papel deja constancia del error; no lo evita.

\vspace{7pt}

% ------------------------------------------------------------
\seccion{La solución --- un sistema ya construido}

\noindent \textbf{PharmaWeigh} controla el surtido y pesaje de materias primas. El operario no avanza si algo no cuadra:

\vspace{6pt}

\noindent
\footnotesize
\begin{tabularx}{\linewidth}{@{}>{\bfseries\color{azulcom}}p{4.0cm} X@{}}
\toprule
Escanea el lote & Con pistola de código de barras o con la cámara del dispositivo. \\
\midrule
El sistema verifica & Que el lote sea del material que pide la receta, que su estado sea \textbf{APROBADO}, que no esté caducado y que alcance la cantidad. Si algo falla, no deja continuar. \\
\midrule
Pesa con semáforo & Verde dentro de tolerancia, ámbar cerca del límite, rojo fuera. Con rojo, el botón de confirmar queda \textbf{deshabilitado}. \\
\midrule
Firma el supervisor & Cada fase de la receta se cierra con el correo y la contraseña de un supervisor o de calidad. El operario no firma su propio trabajo. \\
\midrule
Todo queda asentado & Peso objetivo contra peso real, lote usado, quién dispensó, quién firmó y a qué hora, en una bitácora consultable. \\
\bottomrule
\end{tabularx}
\normalsize

\vspace{6pt}

\begin{cajaInfo}[Ya funciona --- no hay desarrollo por delante]
El sistema está \textbf{construido y desplegado}, con siete roles diferenciados (admin, supervisor, operario, desarrollo, calidad, almacén y auditor), inventario con estados de lote y caducidad, recetas por fases con tolerancias, órdenes de producción, dispensado guiado, firma por fase, bitácora, generación de etiquetas de código de barras, respaldo diario y manuales por rol. Se puede revisar pantalla por pantalla hoy en \href{https://\demoURL}{\demoURL}.
\end{cajaInfo}

\vspace{7pt}

% ------------------------------------------------------------
\needspace{14\baselineskip}
\seccion{Lo que este sistema no es}

\noindent Conviene decirlo antes que después:

\vspace{4pt}

\noindent
\footnotesize
\begin{tabularx}{\linewidth}{@{}L{4.6cm} X@{}}
\toprule
\textbf{No está validado} & No hay CSV ni protocolos IQ/OQ/PQ, ni certificación 21 CFR Part 11 o NOM-059. El sistema tiene funciones \textbf{orientadas} a ese control --- firma por usuario, trazabilidad, bitácora, control de caducidad y de estado de lote --- pero la validación es un proyecto aparte y se cotiza por separado. \\
\midrule
\textbf{No lee la báscula} & El peso se captura a mano en la pantalla. La lectura automática desde báscula está cotizada como opcional. \\
\midrule
\textbf{No incluye hardware} & Tabletas, lectores, básculas e impresoras de etiquetas los pone \clienteCorto. \\
\bottomrule
\end{tabularx}
\normalsize

\vspace{7pt}

% ------------------------------------------------------------
\needspace{14\baselineskip}
\seccion{Lo que falta es de su lado}

\noindent No hay pendientes de desarrollo. Lo que falta para operar son \textbf{decisiones y datos de \clienteCorto}:

\vspace{4pt}

\noindent
\footnotesize
\begin{tabularx}{\linewidth}{@{}X X@{}}
\toprule
\textbullet{} Catálogo de materias primas & \textbullet{} Lotes en piso con caducidad y estado \\
\textbullet{} Formulaciones y \textbf{tolerancias autorizadas} por calidad & \textbullet{} Usuarios, roles y quién queda facultado para firmar \\
\textbullet{} Materiales que se marcan como peligrosos & \textbullet{} Tabletas y lectores por estación de pesaje \\
\bottomrule
\end{tabularx}
\normalsize

\vspace{6pt}

\noindent{\small\color{gristexto}Con eso, MAW Soluciones carga los datos, imprime las etiquetas y ejecuta \textbf{una orden de producción real de punta a punta} en planta --- escaneo, pesaje y firma de todas las fases --- antes de dar por arrancado el sistema. Después acompaña la primera semana.}

\vspace{7pt}

% ------------------------------------------------------------
\needspace{14\baselineskip}
\seccion{Inversión}

\noindent Todo precio sale de la misma fórmula: \textbf{horas de ingeniería $\times$ MXN \tarifa/hora}, en la parte baja del rango de mercado (MXN 900--1,400/h) porque el sistema reutiliza componentes ya probados de MAW. \textbf{Precio cerrado}: lo construido se puede revisar antes de firmar.

\vspace{6pt}

\noindent
\begin{tabularx}{\linewidth}{@{}X C{2.0cm} R{2.8cm} R{3.0cm}@{}}
\toprule
\textbf{Concepto} & \textbf{Horas} & \textbf{Pago único} & \textbf{Mensualidad} \\
\midrule
Sistema construido --- bloques 01 a 10 & \HR{\sTOTAL-(\sK)} & \IM{\sTOTAL-(\sK)} & --- \\
Puesta en marcha con datos reales & \HR{\sK} & \IM{\sK} & --- \\
Operación: servidor, respaldos, soporte y actualizaciones & \HR{\hMa+\hMb+\hMc} & --- & \PENDmini \\
\midrule
\rowcolor{grisfondo}\textbf{PHARMAWEIGH} & \textbf{\HR{\sTOTAL}} & \textbf{\IM{\sTOTAL}} & \textbf{\PENDmini} \\
\bottomrule
\end{tabularx}

\vspace{6pt}

\begin{cajaTotal}
\begin{center}
{\Large\bfseries\color{verdeok}\IM{\sTOTAL} implementación \;$+$\; \PENDIENTE{Fernando}{mensualidad}}\\[4pt]
{\small\color{gristexto}Más IVA. La mensualidad inicia el mes siguiente al arranque. Sin cobro por usuario, por estación, por orden ni por volumen.}
\end{center}
\end{cajaTotal}

\vspace{6pt}

\noindent\textbf{Forma de pago} --- dos pagos, el segundo amarrado a un hecho verificable:

\vspace{4pt}

\noindent
\begin{tabularx}{\linewidth}{@{}C{1.8cm} C{1.2cm} X R{2.6cm}@{}}
\toprule
\textbf{Momento} & \textbf{\%} & \textbf{Libera el pago} & \textbf{Importe} \\
\midrule
Firma & 50\% & Acceso de administrador y puesta en marcha en curso & \IM{(\sTOTAL)*0.5} \\
Arranque & 50\% & Una orden real dispensada y firmada en planta & \IM{(\sTOTAL)*0.5} \\
\bottomrule
\end{tabularx}

\vspace{6pt}

\noindent{\small\color{gristexto}\textbf{Opcionales}, con la misma tarifa: lectura directa desde báscula, validación del sistema (CSV, IQ/OQ/PQ), ampliación de la cobertura de pruebas, integración con ERP, notificaciones por correo y dominio propio. Importes en la sección 5.5 de la propuesta.}

\vspace{7pt}

% ------------------------------------------------------------
\needspace{12\baselineskip}
\seccion{Por qué conviene --- punto de equilibrio}

\noindent El retorno se mide en lotes que no se reprocesan. La tabla se llena con dos datos del cliente, no con estimaciones de MAW:

\vspace{4pt}

\noindent
\begin{tabularx}{\linewidth}{@{}X R{3.6cm} R{3.6cm}@{}}
\toprule
\textbf{Qué se cubre} & \textbf{Lotes reprocesados que se evitan} & \textbf{Costo por orden dispensada} \\
\midrule
La mensualidad & \PENDmini & \PENDmini \\
La implementación & \PENDmini & \PENDmini \\
\bottomrule
\end{tabularx}

\vspace{6pt}

\noindent{\small\color{gristexto}Datos requeridos: \PENDIENTE{cliente}{costo por lote} \quad \PENDIENTE{cliente}{órdenes/mes}}

\vspace{7pt}

% ------------------------------------------------------------
\needspace{10\baselineskip}
\begin{cajaNota}
\textbf{\color{azulcom}Lo que hay que saber antes de firmar}
\begin{itemize}[leftmargin=14pt, itemsep=2pt, topsep=4pt]
    \item Materiales, lotes, formulaciones, órdenes, dispensados, firmas y bitácora son \textbf{propiedad de \clienteCorto}; al terminar el servicio se entrega respaldo completo sin costo.
    \item La \textbf{propiedad intelectual del software es de MAW Soluciones}; el cliente recibe licencia de uso mientras el contrato esté vigente, sin límite de usuarios ni de estaciones.
    \item Alcance adicional se cotiza con la \textbf{misma tarifa} de MXN \tarifa/hora.
    \item El sistema \textbf{no está validado ni certificado}, y \textbf{no lee la báscula}. Ambas cosas son opcionales cotizables.
    \item Las tolerancias y las recetas las autoriza el área de calidad de \clienteCorto; el sistema las aplica, no las define.
\end{itemize}
\end{cajaNota}

\vspace{7pt}

% ------------------------------------------------------------
\needspace{8\baselineskip}
\seccion{Siguiente paso}

\noindent Abrir el demo y recorrer un dispensado completo. Si convence, aprobar la propuesta, firmar el contrato, realizar el \textbf{50\%} y entregar el catálogo de materiales y las formulaciones autorizadas.

\vspace{8pt}

\begin{center}
\begin{tcolorbox}[colback=azulcom, colframe=azulcom, rounded corners, width=13cm, boxrule=1pt, left=6pt, right=6pt, top=6pt, bottom=6pt]
\begin{center}
{\bfseries\color{white}Ver demo --- \href{https://\demoURL}{\color{white}\demoURL}}
\end{center}
\end{tcolorbox}
\end{center}

\vspace{4pt}

\begin{center}
{\small\bfseries Preparado para:} {\bfseries\color{azulcom}\clienteLargo} \quad {\small\bfseries De parte de:} {\bfseries\color{azulcom}MAW Soluciones} \quad {\footnotesize\color{gristexto}\fechaProp}
\end{center}

\end{document}
